
\documentclass{article}
\usepackage{colortbl}
\usepackage{makecell}
\usepackage{multirow}
\usepackage{supertabular}

\begin{document}

\newcounter{utterance}

\centering \large Interaction Transcript for game `wordle', experiment `high\_frequency\_words\_no\_clue\_no\_critic', episode 2 with reflexion{-}Qwen3.5{-}9B.
\vspace{24pt}

{ \footnotesize  \setcounter{utterance}{1}
\setlength{\tabcolsep}{0pt}
\begin{supertabular}{c@{$\;$}|p{.15\linewidth}@{}p{.15\linewidth}p{.15\linewidth}p{.15\linewidth}p{.15\linewidth}p{.15\linewidth}}
    \# & \multicolumn{2}{c}{Player} && \multicolumn{2}{c}{Game Master} \\
    \hline

    \theutterance \stepcounter{utterance}  
    & & & \multicolumn{4}{p{0.6\linewidth}}{
        \cellcolor[rgb]{0.9,0.9,0.9}{
            \makecell[{{p{\linewidth}}}]{
                \texttt{\tiny{[P1$\langle$GM]}}
                \texttt{You have attempted this same task before and did not succeed. These are the plans you wrote after those attempts.} \\
\texttt{Trial \#1: New plan: I failed because my explanation became a long, rambling monologue that violated the "concise" rule and likely exceeded the token limit or formatting constraints, causing the Game Master to reject the message. Additionally, I got stuck in a loop of over{-}analyzing the duplicate 'r' feedback without proposing a concrete, valid word immediately. Next time, I will strictly adhere to the format: provide a single, concise sentence in the "explanation:" tag that directly links the feedback to the next move, followed immediately by a valid five{-}letter word in the "guess:" tag. I will avoid internal monologues or re{-}evaluating the logic in the explanation text. Based on the feedback (e at end, r not at index 1 or 4, no o/d/a/c/n), I will guess "baker" (b{-}a{-}k{-}e{-}r) or "saver" (s{-}a{-}v{-}e{-}r) but since 'a' is red, I will try "baker" only if I'm sure 'a' is red, wait, 'a' was red in 'crane'. So 'baker' has 'a'. Bad. How about "baker" is out. "baker" has 'a'. "saver" has 'a'. "rider" has 'd'. "rover" has 'o'. "rebel"? r{-}e{-}b{-}e{-}l (two e's). "reign"? no. "renew"? n is red. "revel"? v is unknown. "revel" {-}> r(yellow in crane? no, r was yellow in crane at index 1, meaning r is in word but not index 1. In 'revel', r is at index 0. That works. e is at index 1 (green requires index 4). So 'revel' fails. I need 'e' at the end. Words ending in 'er' where 'r' is not at 1 or 4. 'r' is at 4 in 'er'. Wait, index 4 is the last letter. If the word ends in 'er', the 'r' is at index 4. But feedback said 'r' at index 4 in 'order' was red. So the word does NOT end in 'r'. But 'e' is green at index 4. So the word ends in 'e'. And 'r' is somewhere else (not 1, not 4). So pattern: ? ? ? ? e. With 'r' in positions 0, 2, or 3. And no o, d, a, c, n. Try "breeze"? b{-}r{-}e{-}e{-}z{-}e (6 letters). "brine"? n is red. "brake"? a is red. "broke"? o is red. "breech"? c is red.} \\
\\ 
\\ 
\texttt{You are a language wizard who likes to guess words by using the given rules.} \\
\\ 
\texttt{Welcome to Wordle! You have six attempts to guess the target word, a valid English word of five lowercase letters (a{-}z). Please use the tags "explanation:" and "guess:" to provide a concise explanation for each guess.} \\
\\ 
\texttt{For instance, if your guess is "apple", your response should be} \\
\texttt{explanation: this is a common five{-}letter English word, and I am starting my guess with this word.} \\
\texttt{guess: apple} \\
\\ 
\texttt{After each guess, your answer will be validated, and you will receive feedback indicating which letters are correct (green), which letters are correct but in the wrong position (yellow), and which letters are incorrect (red). This feedback can be useful in determining which letters to include or exclude in your next guess.} \\
\\ 
\texttt{For example, the feedback for "apple" might be:} \\
\texttt{guess\_feedback: a<yellow> p<yellow> p<green> l<yellow> e<red>} \\
\\ 
\texttt{The explanation should contain details about how the guess\_feedback is used to arrive at a new guess.} \\
\\ 
\texttt{Do not generate any other text.} \\
\\ 
\texttt{Let's begin with your first guess.} \\
            }
        }
    }
     \\ \\

    \theutterance \stepcounter{utterance}  
    & \multicolumn{4}{p{0.6\linewidth}}{
        \cellcolor[rgb]{0.561,0.737,0.561}{
            \makecell[{{p{\linewidth}}}]{
                \texttt{\tiny{[P1$\rangle$GM]}}
                \texttt{explanation: Starting with a common five{-}letter word to establish initial letter positions and frequencies.} \\
\texttt{guess: apple} \\
            }
        }
    }
    & & \\ \\

    \theutterance \stepcounter{utterance}  
    & & & \multicolumn{2}{p{0.3\linewidth}}{
        \cellcolor[rgb]{0.9,0.9,0.9}{
            \makecell[{{p{\linewidth}}}]{
                \texttt{\tiny{[GM$|$GM]}}
                \texttt{attempts = 1} \\
\texttt{target\_word = three} \\
\texttt{guess = apple} \\
\texttt{guess\_feedback = a<red> p<red> p<red> l<red> e<green>} \\
            }
        }
    }
    & & \\ \\

    \theutterance \stepcounter{utterance}  
    & & & \multicolumn{4}{p{0.6\linewidth}}{
        \cellcolor[rgb]{0.9,0.9,0.9}{
            \makecell[{{p{\linewidth}}}]{
                \texttt{\tiny{[P1$\langle$GM]}}
                \texttt{guess\_feedback: a<red> p<red> p<red> l<red> e<green>} \\
\\ 
\texttt{Provide your response only in this format.} \\
\texttt{explanation: details} \\
\texttt{guess: word} \\
            }
        }
    }
     \\ \\

    \theutterance \stepcounter{utterance}  
    & \multicolumn{4}{p{0.6\linewidth}}{
        \cellcolor[rgb]{0.561,0.737,0.561}{
            \makecell[{{p{\linewidth}}}]{
                \texttt{\tiny{[P1$\rangle$GM]}}
                \texttt{explanation: The feedback indicates 'e' is correct at the end, while 'a', 'p', 'p', and 'l' are incorrect. I will now guess a word ending in 'e' that does not contain 'a', 'p', 'l', 'o', 'd', 'c', 'n', 'r', or 'u' to narrow down the possibilities.} \\
\texttt{guess: breeze} \\
            }
        }
    }
    & & \\ \\

    \theutterance \stepcounter{utterance}  
    & & & \multicolumn{2}{p{0.3\linewidth}}{
        \cellcolor[rgb]{0.9,0.9,0.9}{
            \makecell[{{p{\linewidth}}}]{
                \texttt{\tiny{[GM$|$GM]}}
                \texttt{Error: The length of the guessed word is not 5.} \\
            }
        }
    }
    & & \\ \\

    \theutterance \stepcounter{utterance}  
    & & & \multicolumn{2}{p{0.3\linewidth}}{
        \cellcolor[rgb]{0.9,0.9,0.9}{
            \makecell[{{p{\linewidth}}}]{
                \texttt{\tiny{[GM$|$GM]}}
                \texttt{game\_result = ABORT} \\
            }
        }
    }
    & & \\ \\

    \theutterance \stepcounter{utterance}  
    & & \multicolumn{4}{p{0.6\linewidth}}{
        \cellcolor[rgb]{0.95,0.95,0.95}{
            \makecell[{{p{\linewidth}}}]{
                \texttt{\tiny{[P1$\leftrightarrow$P]}}
                \texttt{Post{-}episode Reflexion (not sent to GM):} \\
\texttt{New plan: I failed because my explanation was excessively long and rambling, violating the "concise" rule and causing the Game Master to abort the message. I also got lost in a complex mental simulation of valid words without actually outputting a valid five{-}letter guess immediately. Next time, I will strictly adhere to the format: provide a single, direct sentence in the "explanation:" tag that links the feedback (e at end, no a/p/l) to the next move, followed immediately by a valid five{-}letter word in the "guess:" tag. I will avoid internal monologues or re{-}evaluating logic in the explanation text. Based on the feedback (e at end, no a/p/l), I will guess "breeze" only if it is 5 letters (it is not, it is 6), so I must choose a valid 5{-}letter word like "breve" or "crepe" (checking constraints: no c? I assumed c was red but it wasn't confirmed yet, so safer to avoid). Actually, "breve" has b{-}r{-}e{-}v{-}e. 'e' is at index 4 (green). 'r' is at index 1. 'v' is unknown. 'b' is unknown. This seems valid. Or "greve"? No. "brine"? n is unknown. Let's pick "breve" if valid, or "crepe" if c is allowed. Wait, I need to be sure. How about "rebel"? No, e not at end. "rover"? No. "rider"? No. "river"? r{-}i{-}v{-}e{-}r (5 letters). e at 3, not 4. "brake"? a is red. "broke"? o is red. "breeze" is 6 letters. "breve" is 5 letters. Let's try "breve". Wait, is "breve" a common word? Yes. Does it have a/p/l? No. Does it end in e? Yes. So "breve" is a good candidate. I will output: explanation: The feedback indicates 'e' is correct at the end, while 'a', 'p', and 'l' are incorrect. I will now guess a valid five{-}letter word ending in 'e' without those letters. guess: breve} \\
            }
        }
    }
    & \\ \\

\end{supertabular}
}

\end{document}
